\documentclass[12pt]{letter}

\usepackage{amsmath}
\usepackage{amssymb}

\usepackage{nopageno}
%\usepackage{amsmath}
\usepackage{enumitem}
% \usepackage{mathtools}
\usepackage{eqnarray,amsmath}
\usepackage[a4paper, total={7.5in, 11.25in}]{geometry}

\begin{document}

Name: \hrulefill \ Neptun code: \hrulefill

\begin{center}
(KTXFI2EBNF) Physics II. Test~2 (Midterm Retake) \\
December~8,~2025 
\end{center}

\bigskip

\begin{enumerate}
\setcounter{enumi}{0} % Számláló beállítása

\item An electron has a kinetic energy of 0.5\,keV.
    \begin{enumerate}
        \item What is its velocity?
        \item What is its momentum?
        \item What is the wavelength of the associated de Broglie wave? \hfill (10~points)
    \end{enumerate}

% ------------------------------------------------------------------------------------------------------------------------------------

\bigskip
\item A metal surface is illuminated with light of wavelength $9 \cdot 10^{-8}\,\text{m}$.  
  What is the velocity of the emitted electrons if the photoelectric effect starts at a wavelength of $2.88 \cdot 10^{-7}\,\text{m}$? \\ $(m_e = 9.1 \cdot 10^{-31}\,\text{kg})$  \hfill (10~points)

% ------------------------------------------------------------------------------------------------------------------------------------

\bigskip
\item In a Compton scattering experiment, X-rays with a wavelength of $0.2\,\text{nm}$ are used. \\ $(m_e = 9.1 \cdot 10^{-31}\,\text{kg})$
    \begin{enumerate}
        \item At what scattering angle does the wavelength of the radiation increase by 1\%?
        \item At what angle does the wavelength become 0.01\% larger? \hfill (10~points)
    \end{enumerate}

% ------------------------------------------------------------------------------------------------------------------------------------

\bigskip
\item A light source has a power output of $20\,\text{W}$.
\begin{enumerate}
    \item What is the frequency of the light if the entire energy of the source is emitted as light with wavelength $\lambda = 800\,\text{nm}$?
    \item How many photons are emitted by the light source per second? \hfill (10~points)
\end{enumerate}

% ------------------------------------------------------------------------------------------------------------------------------------

\bigskip
\item We study an electron in a radioactive particle beam as it passes through the laboratory. It is found that, as measured in the laboratory, the average lifetime of the particle is 60\,ns; after this time the particle decays. The average lifetime measured in the reference frame attached to the particles is $5\,\text{ns}$. What is the propagation velocity of the particle beam? Determine the value of $\frac{m}{m_0}$ for the electron! Calculate the rest energy of an electron corresponding to its rest mass $(m_{e}9.1 \cdot 10^{-31}\,\text{kg})$ in joules and in electronvolts! \\ $(1\,\text{MeV} = 1.6 \cdot 10^{-13}\,\text{J})$ \hfill (10~points)
\end{enumerate}

% ------------------------------------------------------------------------------------------------------------------------------------

Condition for obtaining a signature: achieving at least 40\% (20\,points) on the midterm exam.

Grades: 0–25 points: Fail (1), 26–30 points: Pass (2), 31–37 points: Satisfactory (3), 38–45 points: Good (4), 46–50 points: Excellent (5).

\rightline{Dr.~Gabor FACSKO}
\rightline{\textit{facsko.gabor@uni-obuda.hu}}

\vfill

\end{document}

